\subsection{Artículo especializado para publicar en revista indizada}
Artículo original e inédito, enviado a una revista inscrita en índices validados. (Art. 16, REP).

\subsubsection*{Requisitos}
\begin{itemize}
  \item Artículo original e inédito, con autoría del pasante.
  \item Tema relacionado con el plan de estudios o con un área científica o técnica de la profesión.
  \item Evidencia de que la revista está en alguno de los índices del Art. 16 (e-Revist@s, Redalyc, Scielo, Scopus, CONACYT, Thomson-Reuters, JCR).
  \item Enviar el artículo hasta contar con el voto aprobatorio del asesor.
  \item Reporte de similitud iThenticate no mayor al 20 \%.
\end{itemize}
\capturapendiente{Pantalla de opciones de titulación.}
\capturapendiente{Requisitos de la opción.}
\subsubsection*{Etapa 1. Tus datos}
Captura una sola vez tus datos; con ellos se llenan automáticamente todos los formatos del trámite.

Datos que se capturan:
\begin{itemize}
  \item Nombre completo (con acentos, tal como aparecerá en el título)
  \item Número de cuenta
  \item Licenciatura
  \item Sistema
  \item Plan de estudios (clave y versión)
  \item Escuela incorporada a la UAEM donde cursaste la licenciatura
  \item Año de ingreso
  \item Año de egreso
  \item Teléfono
  \item Correo institucional
  \item Correo personal
  \item Ocupación actual
  \item Fecha de nacimiento
  \item Mi nombre o apellidos llevan acentos
\end{itemize}
\capturapendiente{Captura de datos del pasante.}
\subsubsection*{Etapa 2. Protocolo y aprobación del tema}
Captura la propuesta de tu trabajo. El sistema llena el formato oficial de solicitud y protocolo; escribe aquí todo el protocolo; la plataforma lo firma por ti, le pide la firma a tu asesor por correo y lo envía al responsable del Comité de Revisión de Protocolos de tu licenciatura. Si el comité te pide correcciones, las verás en tus Notificaciones: corrige aquí y vuelve a enviarlo.

Datos que se capturan:
\begin{itemize}
  \item Título tentativo del artículo
  \item Área académica
  \item Asesor (adscrito a la UAEM): nombre
  \item Asesor: correo electrónico
  \item Asesor: grado (Dr., Mtra., Ing.…)
  \item Asesor: género
  \item Coasesor (opcional): nombre
  \item Coasesor: correo electrónico
  \item Coasesor: grado
  \item Coasesor: género
  \item Nombre de la revista indizada
  \item Índice al que pertenece (nombre y URL)
  \item Descripción general (resumen) del artículo
  \item Referencias y/o fuentes de información
\end{itemize}

La plataforma llena los formatos: protocolo y los firma por ti.

Documentos (PDF):
\begin{longtable}{p{0.55\linewidth}p{0.35\linewidth}}
\toprule
\textbf{Documento} & \textbf{Observaciones}\\
\midrule
\endhead
Solicitud y protocolo firmados (pasante y asesor) & La plataforma lo genera. La plataforma lo genera con lo que escribiste; tu asesor lo firma por correo. Después lo revisa el comité de tu licenciatura.\\
Carta de presentación del pasante y asesor & Descárgala en blanco desde «Formatos en blanco». Solo para instituciones incorporadas.\\
\bottomrule
\end{longtable}

\emph{Turno del departamento:} Aprobación del protocolo y oficio de aprobación del tema. Cuando el comité aprueba el protocolo, la plataforma genera el oficio de aprobación del tema; lo firman la Jefatura de Evaluación Profesional (Vo. Bo.) y la Subdirección Académica.
\capturapendiente{Etapa en curso del trámite.}
\subsubsection*{Etapa 3. Desarrollo del trabajo escrito}
Tienes exclusividad sobre el tema durante dos años a partir de la aprobación; puedes renovarla con una prórroga hasta por un año, que incluye la sustentación (Art. 86, REP). Desarróllalo con tu asesor; cuando cuente con su visto bueno, confirma aquí para continuar.

Al terminar, confirma: «Concluí el trabajo escrito y mi asesor lo avala para revisión».
\subsubsection*{Etapa 4. Solicitud de revisión del trabajo escrito (examen previo)}
Llena la solicitud de revisión del trabajo escrito y sube los documentos probatorios: el reporte de similitud de iThenticate (máximo 20 \%) y tu trabajo en PDF. Después elige la fecha y hora de tu examen previo y propone a tu sínodo; tu solicitud llega al Coordinador de Docencia de tu licenciatura.

Datos que se capturan:
\begin{itemize}
  \item Título del trabajo escrito
\end{itemize}

La plataforma llena los formatos: revision y los firma por ti.

Documentos (PDF):
\begin{longtable}{p{0.55\linewidth}p{0.35\linewidth}}
\toprule
\textbf{Documento} & \textbf{Observaciones}\\
\midrule
\endhead
Solicitud de revisión firmada (pasante y asesor) & La plataforma lo genera. La firmas tú y tu asesor da su Vo. Bo. desde su correo.\\
Reporte de porcentaje de similitud (iThenticate) & No mayor al 20 \%.\\
Trabajo escrito en PDF para revisión & Con copia del oficio de aprobación de tema al inicio.\\
Solicitud de prórroga del tema (pasante y asesor) & La plataforma lo genera. Tu protocolo se aprobó hace más del plazo de exclusividad (Art. 86, REP). Captura avance y motivos en «Solicitudes y formatos adicionales», firma aquí y tu asesor da su Vo. Bo.; al validarla, la plataforma genera el oficio de aprobación de prórroga que firma la Subdirección.\\
\bottomrule
\end{longtable}

Elige fecha y hora, propone a tu sínodo o jurado y envía la solicitud al Coordinador de Docencia.

\emph{Turno del departamento:} Fecha del examen previo y sínodo. Eliges fecha y hora y propones tu sínodo; el Coordinador de Docencia de tu licenciatura lo ratifica o modifica y firma desde el enlace de su correo. Después se generan los oficios individuales.
\subsubsection*{Etapa 5. Solicitud de sustentación y voto aprobatorio}
Cuando hayas atendido las correcciones del examen previo, genera el voto aprobatorio: la plataforma lo envía a tu asesor y a tus revisores para que lo firmen desde su correo. Con el voto firmado, firma tu solicitud de sustentación (tu asesor da su Vo. Bo.) para obtener la autorización de impresión y la asignación del jurado. Escribe aquí la estructura capitular: con ella la plataforma genera el oficio de estructura capitular que firma el Director.

Datos que se capturan:
\begin{itemize}
  \item Estructura capitular (un capítulo por línea, sin numerar)
\end{itemize}

La plataforma llena los formatos: sustentacion y los firma por ti.

Documentos (PDF):
\begin{longtable}{p{0.55\linewidth}p{0.35\linewidth}}
\toprule
\textbf{Documento} & \textbf{Observaciones}\\
\midrule
\endhead
Voto aprobatorio (lo firman tu asesor y tus revisores) & La plataforma lo genera. La plataforma lo genera y pide las firmas por correo.\\
Solicitud de sustentación firmada (pasante y asesor) & La plataforma lo genera. Se habilita cuando el voto aprobatorio está firmado.\\
\bottomrule
\end{longtable}

\emph{Turno del departamento:} Autorización de impresión. 
\subsubsection*{Etapa 6. Documentación para la evaluación profesional}
Reúne y sube cada documento en PDF. Cuando todos estén validados, Apoyo a Titulación te cita para entregar originales, copias y ejemplares, y fijar el jurado.

La plataforma llena los formatos: trabajo\_eval, datos\_titulo, requisitos y los firma por ti.

Documentos (PDF):
\begin{longtable}{p{0.55\linewidth}p{0.35\linewidth}}
\toprule
\textbf{Documento} & \textbf{Observaciones}\\
\midrule
\endhead
Solicitud para presentar la evaluación profesional (trabajo escrito) & La plataforma lo genera. La plataforma lo genera con tus datos; lo firmas electrónicamente y la Subdirección lo autoriza.\\
Trabajo escrito final en PDF & Con oficio de autorización de impresión, resumen y carta de autorización para el Repositorio Institucional.\\
Reporte de porcentaje de similitud (iThenticate) & No mayor al 20 \%.\\
Constancia de pertenencia de la revista a los índices validados & Emitida por la Secretaría de Investigación y Estudios Avanzados (Art. 16, REP).\\
Carta de recepción o publicación del artículo expedida por la revista & Incluye criterios de publicación y documento de envío (Art. 20, REP).\\
Certificado de estudios de licenciatura & Descargado del Portal de alumnos.\\
Certificado de servicio social & \\
Certificado de estudios de bachillerato & \\
Constancia de autenticación de bachillerato & Se tramita en la escuela donde cursaste el bachillerato. Solo si tus estudios de bachillerato no fueron en instituciones dependientes o incorporadas a la UAEM.\\
Acta de nacimiento & Copia certificada vigente (del año en curso).\\
CURP & Debe mostrar la leyenda «Validada por el Registro Civil».\\
Constancia o certificado de aprobación del idioma inglés & Solo planes 98 en adelante, si en tu certificado no aparecen Inglés C1 e Inglés C2.\\
Constancia de no adeudo a la Biblioteca de la Facultad & El recibo de pago se tramita en Control Escolar de la Facultad; la constancia, en la Biblioteca.\\
Constancia de no adeudo a la Biblioteca Central & \\
Contra-recibo de la donación de libros a la Biblioteca de la Facultad & \\
Constancia de no adeudo a la Dirección de Control Escolar de la UAEM & Se solicita en el Portal de alumnos: controlescolar.uaemex.mx.\\
Recibo de no adeudo a la UAEM & Se solicita en el Portal de alumnos de Control Escolar.\\
Formato de datos personales (para el título) & La plataforma lo genera. La plataforma lo llena con tus datos; lo firmas electrónicamente.\\
Comprobante del sistema de Seguimiento de Egresados & Se obtiene en alumniversitario.uaemex.mx/1dae.php.\\
Dictamen de autorización de los HH. Consejos & La plataforma lo genera. Se solicita en la Subdirección Académica con el formato del Art. 5. Solo si excediste el plazo de 10 años para presentar la evaluación profesional (Art. 5, REP).\\
Carta de autorización o de excepción para publicar en el Repositorio Institucional & Entrega la que corresponda; descárgala en «Formatos en blanco».\\
Recibo de pago: derechos de evaluación profesional (\$3,024) & Se solicita en Control Escolar de la Facultad.\\
Recibo de pago: expedición de título (\$1,187) & \\
\bottomrule
\end{longtable}
\subsubsection*{Etapa 7. Solicitud del examen definitivo}
Elige la fecha y hora de tu examen definitivo y propone a tu jurado; tu solicitud llega al Coordinador de Docencia de tu licenciatura, que lo ratifica o modifica.

Elige fecha y hora, propone a tu sínodo o jurado y envía la solicitud al Coordinador de Docencia.

\emph{Turno del departamento:} Fecha del examen definitivo y jurado. Eliges fecha y hora y propones tu jurado; el Coordinador de Docencia lo ratifica o modifica y firma desde el enlace de su correo. Después se generan los oficios individuales.
\subsubsection*{Etapa 8. Resultado de la evaluación}
Al concluir, el departamento registra el resultado. Después podrás consultar cómo descargar tu título y cédula profesional.

\emph{Turno del departamento:} Registrar resultado. 

\subsubsection*{Seguimiento y conclusión}
\capturapendiente{Trámite concluido.}
\capturapendiente{Ruta del trámite.}
\aviso{Esta opción aún no se ha recorrido en la prueba de extremo a extremo con capturas; el texto sale del catálogo del sistema.}
